% ------------------------------------------------------------------------------------------
% The Canonical Cryptography Bill of Materials (2CBOM) -- Formal Specification, version 1.0
% Compiles with pdfLaTeX (Overleaf default) using only standard TeX Live packages.
% ------------------------------------------------------------------------------------------
\documentclass[11pt,a4paper]{article}
\usepackage[utf8]{inputenc}
\usepackage[T1]{fontenc}
\IfFileExists{lmodern.sty}{\usepackage{lmodern}}{}
\IfFileExists{lmodern.sty}{\usepackage{microtype}}{\usepackage[expansion=false]{microtype}}
\usepackage[margin=2.4cm]{geometry}
\usepackage{amsmath,amssymb,amsthm}
\usepackage{booktabs,longtable,array,tabularx}
\usepackage{enumitem}
\usepackage{xcolor}
\usepackage{listings}
\PassOptionsToPackage{hyphens,obeyspaces,spaces}{url}
\usepackage[hidelinks,pdftitle={2CBOM 1.0 Formal Specification},pdfauthor={Amin Jafarpour}]{hyperref}
\emergencystretch=2.5em

% ---------- typography ----------
% \code{...}: monospace that may break after - / . : and spaces (text mode);
% \codem{...}: unbreakable monospace for use inside mathematics.
\DeclareUrlCommand\code{\urlstyle{tt}}
\newcommand{\codem}[1]{\texttt{\detokenize{#1}}}
\newcommand{\twocbom}{2CBOM}
\newcommand{\MUST}{\textbf{MUST}}
\newcommand{\MUSTNOT}{\textbf{MUST NOT}}
\newcommand{\SHOULD}{\textbf{SHOULD}}
\newcommand{\MAY}{\textbf{MAY}}
\newcommand{\sem}[1]{\llbracket #1 \rrbracket}
\newcommand{\llbracket}{[\![}
\newcommand{\rrbracket}{]\!]}
\newcommand{\pfun}{\rightharpoonup}
\newcommand{\Ids}{\mathcal{K}}
\newcommand{\Assets}{\mathcal{A}}
\newcommand{\Comps}{\mathcal{C}}
\newcommand{\Rels}{\mathcal{R}}
\newcommand{\Evs}{\mathcal{E}}
\newcommand{\Srcs}{\mathcal{S}}
\newcommand{\Finds}{\mathcal{F}}
\newcommand{\Facts}{\mathcal{W}}
\newcommand{\Q}{\mathcal{Q}}
\newcommand{\lfp}{\operatorname{lfp}}
\newcommand{\worst}{\operatorname{worst}}
\newcommand{\canon}{\operatorname{canon}}
\newcommand{\id}{\operatorname{id}}
\newcommand{\digest}{\operatorname{digest}}
\newcommand{\JCS}{\operatorname{JCS}}
\newcommand{\merge}{\oplus}

\lstset{basicstyle=\ttfamily\footnotesize,columns=fullflexible,keepspaces=true,breaklines=true,
        frame=single,framerule=0.3pt,xleftmargin=4pt,xrightmargin=4pt,showstringspaces=false}

\theoremstyle{definition}
\newtheorem{definition}{Definition}[section]
\newtheorem{example}[definition]{Example}
\newtheorem{remark}[definition]{Remark}
\theoremstyle{plain}
\newtheorem{theorem}{Theorem}
\newtheorem{lemma}[theorem]{Lemma}
\newtheorem{proposition}[theorem]{Proposition}
\makeatletter\newtheorem{constraint}{Constraint}
\renewcommand{\theconstraint}{C\two@digits{\value{constraint}}}\makeatother

\title{\textbf{The Canonical Cryptography Bill of Materials (\twocbom)}\\[4pt]
       \large A Formal Specification for Post-Quantum Cryptography Compliance\\[2pt]
       \normalsize Version 1.0}
\author{Amin Jafarpour\\ School of Computing Science, Simon Fraser University}
\date{October 29, 2026}

\begin{document}
\maketitle

\begin{abstract}
Cryptography Bill of Materials (CBOM) generators disagree on names, structure, completeness
and assessment of the cryptographic assets they report, and the pre-standard IBM CBOM format
coexists with CycloneDX~1.6 and~1.7 documents. This specification defines \twocbom, a canonical
model for cryptographic inventories that is strict enough to support post-quantum cryptography
(PQC) compliance determinations. \twocbom{} fixes (i) a deterministic canonical identity for
algorithms, certificates, key material, protocols and software components, (ii) closed
vocabularies, (iii) structural closure constraints that make every asset assessable,
(iv) a mandatory quantum-security assessment per asset with NIST~IR~8547, CNSA~2.0 and
FIPS status, (v) a canonical serialisation with content digest and deterministic serial
number, and (vi) merge semantics under which the reconciliation of any set of source CBOMs
is a commutative, associative and idempotent operation whose result is independent of the
order in which the sources are processed. Mappings from CycloneDX~1.5/1.6/1.7 and from the
IBM \code{1.4-cbom-1.0} format are given, together with a lossless-by-annotation projection
back to CycloneDX~1.6. A reference implementation, \code{twocbom}, realises the model as a
rule-based expert system.
\end{abstract}

\tableofcontents

% ==========================================================================================
\section{Introduction and scope}
% ==========================================================================================
\subsection{Motivation}
Regulatory and organisational PQC programmes (for instance the NIST transition timeline of
IR~8547, which deprecates 112-bit quantum-vulnerable public-key cryptography after 2030 and
disallows all quantum-vulnerable public-key cryptography after 2035, and the CNSA~2.0 suite)
require an \emph{auditable inventory} of cryptographic assets. CBOM generators produce such
inventories, but in practice the outputs of CBOMkit, cdxgen, cryptobom-forge, cbomkit-theia and
earlier IBM tooling differ in at least four ways:
\begin{enumerate}[nosep]
  \item \emph{Naming}: \code{AES128-GCM}, \code{AES/GCM/NoPadding}, \code{aes-128-gcm} and the OID
        \code{2.16.840.1.101.3.4.1.6} denote the same algorithm.
  \item \emph{Structure and version}: CycloneDX~1.6 \code{cryptographic-asset} components with
        \code{evidence.occurrences} versus IBM \code{crypto-asset} components with
        \code{detectionContext} and camelCase enumerations.
  \item \emph{Completeness}: one tool sees the key size, another the mode of operation, a third only
        the family.
  \item \emph{Assessment}: security levels copied from different tables, or absent.
\end{enumerate}
A compliance decision taken over such inputs is neither reproducible nor explainable.

\subsection{Scope}
\twocbom{} specifies (a) a data model and closed vocabularies; (b) canonical identity;
(c) well-formedness and conformance constraints in two profiles, \emph{core} and \emph{strict};
(d) a quantum-security assessment model; (e) a canonical serialisation; (f) merge semantics;
(g) mappings from and to existing formats. It does not specify how assets are detected.

\subsection{Conventions}
The key words \MUST, \MUSTNOT, \SHOULD{} and \MAY{} are to be interpreted as in RFC~2119.
$\Sigma^{*}$ denotes finite strings over the Unicode alphabet; $\bot$ denotes an absent value;
$f: X \pfun Y$ denotes a partial function; $\mathcal{P}(X)$ the power set of $X$;
$\mathcal{P}^{+}(X)$ the non-empty subsets. JSON member names are written in \code{monospace}.
Identifiers are compared byte-wise after Unicode normalisation form~C.

% ==========================================================================================
\section{Overview of the model}
% ==========================================================================================
A \twocbom{} document is a finite, typed, acyclic-in-identity structure
$D = (\textit{hdr}, \mathcal{L}, \Comps, \Assets, \Rels, \Finds, \Gamma)$ consisting of a header,
\emph{lineage} $\mathcal{L}$ (the sources and the decisions taken while merging them), software
\emph{components} $\Comps$, cryptographic \emph{assets} $\Assets$, typed \emph{relations} $\Rels$
between components and assets, \emph{findings} $\Finds$ produced by conformance checking, and a
\emph{conformance} record $\Gamma$. Every entity carries a canonical identifier from the
identifier space $\Ids$ (Section~\ref{sec:identity}); every asset carries evidence, a set of
contributing sources and an assessment (Section~\ref{sec:assessment}). The document is
serialised canonically and sealed with a digest (Section~\ref{sec:serialisation}). Section~\ref{sec:merge}
defines how a set of source CBOMs is turned into a \twocbom{} document.

% ==========================================================================================
\section{Core data model}
% ==========================================================================================
\subsection{Vocabularies}
\begin{definition}[Closed vocabularies]\label{def:vocab}
The following finite sets are \emph{closed}: a document is non-conformant if an enumerated
attribute carries a value outside the corresponding set. Each set is a superset of the
CycloneDX~1.6 enumeration of the same purpose, so that every value has a CycloneDX image
(Section~\ref{sec:projection}).
\begin{description}[nosep,leftmargin=2.6cm,style=nextline]
\item[$V_{\mathrm{type}}$] \{algorithm, certificate, related-crypto-material, protocol\}
\item[$V_{\mathrm{prim}}$] \{drbg, mac, block-cipher, stream-cipher, signature, hash, pke, xof, kdf, key-agree, kem, ae, combiner, other, unknown\}
\item[$V_{\mathrm{fn}}$] \{generate, keygen, encrypt, decrypt, digest, tag, keyderive, sign, verify, encapsulate, decapsulate, other, unknown\}
\item[$V_{\mathrm{mode}}$] \{cbc, ecb, ccm, gcm, cfb, ofb, ctr, xts, siv, gcm-siv, ocb, eax, cts, kw, kwp, other, unknown\}
\item[$V_{\mathrm{pad}}$] \{pkcs5, pkcs7, pkcs1v15, oaep, pss, raw, none, zero, iso10126, iso7816, x923, other, unknown\}
\item[$V_{\mathrm{curve}}$] \{p192, p224, p256, p384, p521, secp256k1, brainpoolp256r1, brainpoolp384r1, brainpoolp512r1, curve25519, curve448, sm2p256v1, other, unknown\}
\item[$V_{\mathrm{env}}$] \{software-plain-ram, software-encrypted-ram, software-tee, hardware, other, unknown\}
\item[$V_{\mathrm{cert}}$] \{none, fips140-$k$-l$n$ ($k\in\{1,2,3\}$, $n\in\{1,\dots,4\}$), cc-eal$n$, cc-eal$n$+ ($n\in\{1,\dots,7\}$), other, unknown\}
\item[$V_{\mathrm{proto}}$] \{tls, ssh, ipsec, ike, sstp, wpa, other, unknown\}
\item[$V_{\mathrm{mat}}$] \{private-key, public-key, secret-key, key, ciphertext, signature, digest, initialization-vector, nonce, seed, salt, shared-secret, tag, additional-data, password, credential, token, other, unknown\}
\item[$V_{\mathrm{state}}$] \{pre-activation, active, suspended, deactivated, compromised, destroyed\}
\item[$V_{\mathrm{rel}}$] \{uses, implements, provides, negotiates, signed-by, subject-key-algorithm, instantiates, secured-by, depends-on, component-of\}
\item[$V_{\mathrm{method}}$] \{static-source, static-binary, dynamic, network, configuration, manual, declared, derived, unknown\}
\item[$\Q$] \{classically-weak, quantum-vulnerable, unknown, quantum-safe\} (quantum classes)
\item[$V_{\mathrm{ir}}$] \{acceptable, deprecated-after-2030, disallowed-after-2035, disallowed, unknown\}
\item[$V_{\mathrm{prio}}$] \{immediate, high, medium, low, none, unknown\}
\item[$V_{\mathrm{sev}}$] \{error, warning, info\}
\end{description}
\end{definition}

\subsection{Entities}
\begin{definition}[Source record]
A source record is a tuple $s = (\id_s, \delta_s, \phi_s, v_s, t_s, \tau_s, \mu_s)$ with a
source identifier $\id_s = \codem{src-}h$ where $h$ is the first 12 hexadecimal digits of the
SHA-256 digest $\delta_s$ of the source bytes, a format $\phi_s \in$ \{cyclonedx-json,
ibm-cbom-json, cyclonedx-prestandard-json\}, the source's own spec version $v_s$, a tool
record $t_s = (\text{name}, \text{version})$, a trust value $\tau_s \in \{0,\dots,100\}$ and a
default detection method $\mu_s \in V_{\mathrm{method}}$. $\Srcs$ denotes the set of source
records of a document. Because $\id_s$ is derived from content, the same source file always
receives the same identifier.
\end{definition}

\begin{definition}[Evidence record]\label{def:evidence}
An evidence record is a tuple $e = (\id_s, t, \mu, \ell, n, o, \sigma, \gamma, c)$: the
contributing source and its tool, a detection method $\mu \in V_{\mathrm{method}}$, an optional
normalised location $\ell \in \Sigma^{*} \cup \{\bot\}$ (file path, binary, endpoint), optional
line $n$ and offset $o \in \mathbb{N} \cup \{\bot\}$, optional symbol $\sigma$ and context
$\gamma$, and a confidence $c \in [0,1]$. Evidence with $\mu = \text{derived}$ \MUST{} name in
$\gamma$ the entity and the rule it was derived from. The \emph{key} of $e$ is
$(\id_s,\ell,n,o,\sigma,\mu)$; two evidence records with equal keys are the same record.
\end{definition}

\begin{definition}[Algorithm descriptor]\label{def:descriptor}
Let $\mathbb{F}$ be the set of algorithm \emph{families} of the knowledge base
(\code{aes}, \code{rsa}, \code{ml-kem}, \code{sha2}, \ldots) and
$\mathbb{F}^{+} = \mathbb{F} \cup \{f_1\text{+}\cdots\text{+}f_k : f_i \in \mathbb{F}, k\ge 2\}$
the families extended by \emph{hybrid combiners}. An algorithm descriptor is a tuple
$$\delta = (f,\ \pi,\ \kappa,\ m,\ p,\ h,\ \rho,\ \chi)$$
with family $f \in \mathbb{F}^{+}$, parameter set $\pi \in \Sigma^{*}\cup\{\bot\}$ (key bits,
modulus bits, digest bits, or a PQ parameter-set name such as \code{sha2-128s}), curve
$\kappa \in V_{\mathrm{curve}}\cup\{\bot\}$, mode $m \in V_{\mathrm{mode}}\cup\{\bot\}$, padding
$p \in V_{\mathrm{pad}}\cup\{\bot\}$, hash $h$ (the canonical string of a hash descriptor, or
$\bot$), derived primitive $\rho \in V_{\mathrm{prim}}$ and constituents $\chi$ (a sequence of
canonical strings; non-empty exactly for hybrids and cipher-suite-shaped assets). The knowledge
base assigns to each family the set of \emph{admissible slots}
$\mathrm{slots}(f) \subseteq \{\pi,\kappa,m,p,h\}$; a descriptor \MUST{} leave inadmissible
slots at $\bot$ (\code{RSA/ECB/...} therefore has no mode).
\end{definition}

\begin{definition}[Asset]\label{def:asset}
An asset is a tuple $a = (\kappa_a, \tau_a, \nu_a, \Lambda_a, P_a, \alpha_a, E_a, S_a, u_a)$ with
canonical identifier $\kappa_a \in \Ids$, type $\tau_a \in V_{\mathrm{type}}$, display name
$\nu_a$, alias set $\Lambda_a \subseteq \Sigma^{*}$ (every raw name that was mapped to $\kappa_a$),
a typed property record $P_a$ (Definition~\ref{def:props}), an assessment $\alpha_a$
(Definition~\ref{def:assessment}), a non-empty evidence set $E_a \subseteq \Evs$, the set
$S_a \subseteq \Srcs$ of contributing sources, and an \emph{unresolved} flag $u_a \in \{0,1\}$
that is $1$ exactly when the identity could not be canonicalised.
\end{definition}

\begin{definition}[Property records]\label{def:props}
$P_a$ is determined by $\tau_a$:
\begin{description}[nosep,leftmargin=1.5em]
\item[algorithm] the descriptor slots of $\delta$, the primitive, the crypto functions
  $\subseteq V_{\mathrm{fn}}$, the classical security level $\lambda_c \in \mathbb{N}\cup\{\bot\}$,
  the NIST quantum security level $\lambda_q \in \{0,\dots,5\}\cup\{\bot\}$, the OID, the governing
  standards, execution environment, implementation platform, certification levels, pre-standard
  names, and the constituents $\chi$ (as identifiers).
\item[certificate] subject, issuer, serial number, \code{notBefore}, \code{notAfter} (ISO~8601,
  UTC), format, extension, fingerprint, and the identifiers of the signature algorithm and of the
  subject public-key algorithm (through relations \code{signed-by} and \code{subject-key-algorithm}).
\item[related-crypto-material] type $\in V_{\mathrm{mat}}$, identifier, state, size in bits,
  format, dates, the algorithm(s) it instantiates (relation \code{instantiates}), the mechanism
  securing it (relation \code{secured-by}) and fingerprint. The material \emph{value} itself
  \MUSTNOT{} be present (Section~\ref{sec:security}).
\item[protocol] type $\in V_{\mathrm{proto}}$, normalised version, the set of cipher suites
  (each a tuple of name, identifier, the identifiers of its constituent algorithms and the
  flag \emph{kexKnown}), and references to further assets (relation \code{uses}).
\end{description}
\end{definition}

\begin{definition}[Component]
A software component is $c = (\kappa_c, \tau_c, \nu_c, \mathrm{ver}, \mathrm{grp}, \mathrm{purl},
\mathrm{cpe}, \Lambda_c, S_c, r_c)$ with type $\tau_c$ in the CycloneDX component-type
vocabulary and root flag $r_c$ (the subject of the inventory).
\end{definition}

\begin{definition}[Relation]
A relation is $r = (\kappa_1, \kappa_2, \rho, \text{via}, S_r)$ with $\kappa_1,\kappa_2 \in \Ids$,
kind $\rho \in V_{\mathrm{rel}}$, an optional qualifier \emph{via} (the cipher-suite name for
\code{negotiates}) and contributing sources $S_r$. The key of $r$ is $(\kappa_1,\kappa_2,\rho,\text{via})$.
\end{definition}

\begin{definition}[Finding and conflict]
A finding is $(\text{code}, \text{severity} \in V_{\mathrm{sev}}, \text{subject}, \text{message},
\text{rule})$. A conflict record is $(\text{subject}, \text{attribute}, \text{chosen},
\text{candidates}, \text{rule})$ where \emph{candidates} lists every value asserted for the
attribute together with its origin, strength, trust and confidence.
\end{definition}

\begin{definition}[Document]\label{def:document}
A \twocbom{} document is a tuple
$$D = (\textit{ver}, \textit{profile}, \textit{sn}, \textit{gen}, \mathcal{I}, \mathcal{L},
\Comps, \Assets, \Rels, \Finds, \Gamma, \Sigma_D)$$
where $\textit{ver}=\codem{1.0}$;
$\textit{profile} \in \{\text{core}, \text{strict}\}$; $\textit{sn}$ is the serial number and
$\mathcal{I}$ the integrity block (Section~\ref{sec:serialisation}); $\textit{gen}$ records
generator and timestamp; $\mathcal{L} = (\Srcs, \text{ruleset}, \text{knowledge base},
\text{conflicts})$ is the lineage; $\Gamma = (\textit{profile}, \text{conformant}, \#\text{errors},
\#\text{warnings}, \#\text{infos})$ and $\Sigma_D$ a summary.
\end{definition}

% ==========================================================================================
\section{Canonical identity}\label{sec:identity}
% ==========================================================================================
\subsection{Identifier space}
\begin{definition}[Identifier grammar]
$\Ids \subset \Sigma^{*}$ is the language generated by
\begin{center}\ttfamily\small
\begin{tabular}{l@{\ ::=\ }p{12.2cm}}
id        & "2cbom:" kind ":" body \\
kind      & "alg" | "cert" | "key" | "proto" | "comp" \\
body      & canonical-string \quad(kind = alg)\\
          & \ |\ ("fp:" hashalg ":" hex) | ("dn:" hex) | ("loc:" hex) | ("unresolved:" hex) \quad(kind = cert)\\
          & \ |\ mattype ":" ("fp:" hashalg ":" hex | "id:" hex | "loc:" hex | "src:" hex) \quad(kind = key)\\
          & \ |\ prototype ["-" version] \quad(kind = proto)\\
          & \ |\ purl | cpe | [group "/"] name ["@" version] \quad(kind = comp)
\end{tabular}
\end{center}
where \code{hex} is a lower-case hexadecimal string and \code{hashalg} $\in$ \{sha256, sha1, sha384, sha512, md5\}.
\end{definition}

\subsection{Canonical strings of algorithms}
\begin{definition}[Canonical string]\label{def:canonstring}
For a descriptor $\delta=(f,\pi,\kappa,m,p,h,\rho,\chi)$ the canonical string is the
\code{-}-joined concatenation of the defined slots in the fixed \emph{slot order}
$$\sem{\delta} = f \;\cdot\; \pi \;\cdot\; \kappa \;\cdot\; m \;\cdot\; p \;\cdot\; h,$$
omitting $\bot$ slots, and $\id(\delta) = \codem{2cbom:alg:}\sem{\delta}$. Hybrid families are
written $f_1\text{+}f_2$ where each $f_i$ is itself a canonical string
(\code{x25519+ml-kem-768}). Slot values are drawn from the grammar
\begin{center}\ttfamily\small
\begin{tabular}{l@{\ ::=\ }p{11.2cm}}
canonical-string & family ["-" pset] ["-" curve] ["-" mode] ["-" padding] ["-" hash] \\
family           & token | token "+" family \\
pset             & digits | digits "-" digits | pq-set | variant \\
pq-set           & ("sha2" | "shake") "-" digits ("s" | "f") \\
hash             & "sha1" | "md5" | ("sha2" | "sha3" | "shake" | "blake2b" | ...) "-" digits \\
token            & [a-z0-9]+ ("-" [a-z0-9]+)*
\end{tabular}
\end{center}
\end{definition}
Because the slot order is fixed and each slot is drawn from a closed set, two descriptors are
equal iff their canonical strings are equal (the map $\delta \mapsto \sem{\delta}$ is injective on
descriptors that respect $\mathrm{slots}(f)$).

\begin{example}
\begin{itemize}[nosep]
\item \code{2cbom:alg:aes-128-gcm} (from \code{AES128-GCM}, \code{AES-128-GCM}, \code{aes-128-gcm}, OID \code{2.16.840.1.101.3.4.1.6});
\item \code{2cbom:alg:rsa-2048-oaep-sha2-256}; \code{2cbom:alg:ecdsa-p256-sha2-256}; \code{2cbom:alg:tdea-cbc};
\item \code{2cbom:alg:rsa-pkcs1v15-sha2-256} (the Java name \code{SHA256withRSA}, the X.509 name
      \code{sha256WithRSAEncryption} and the OID \code{1.2.840.113549.1.1.11});
\item \code{2cbom:alg:ml-kem-768} (also \code{Kyber768}); \code{2cbom:alg:slh-dsa-shake-256f}; \code{2cbom:alg:hmac-sha2-256};
\item \code{2cbom:alg:x25519+ml-kem-768} (hybrid combiner).
\end{itemize}
\end{example}

\subsection{The canonicalisation function}
\begin{definition}[Canonicalisation]\label{def:canon}
Let $H$ be the set of \emph{hint records} (partial maps from \{primitive, parameterSetIdentifier,
curve, mode, padding, oid\} to strings, as carried by structured CBOM fields). The
canonicalisation function
$$\canon : \Sigma^{*} \times H \to (\Delta \times [0,1] \times \mathcal{P}(\Sigma^{*})) \cup \{\bot\}$$
maps a raw name and hints to a descriptor, a confidence and a set of notes, or to $\bot$
(\emph{unresolved}). It is defined by the following stages, each of which is deterministic:
\begin{enumerate}[nosep]
\item \textbf{OID stage.} If the name (or the hint \code{oid}) is an OID present in the knowledge
      base, the descriptor registered for the OID is used; a parsable name may additionally
      refine slots the OID leaves open.
\item \textbf{Composite stage.} Names of the form \code{TLS_...} or \code{SSL_...} are cipher
      suites and yield $f = \codem{suite}$ with constituents from Definition~\ref{def:suite}.
\item \textbf{Alias stage.} The name, lower-cased and stripped of separators
      (hyphen, underscore, slash, space, comma, dot, plus, brackets, colon, semicolon, quotes), is looked up in the whole-name lexicon
      (\code{x25519mlkem768}, \code{rsassapss}, \code{desede}, \code{securerandom}, \ldots).
\item \textbf{Segmentation stage.} Otherwise the squashed name is segmented by greedy
      longest-match over the lexicon of \emph{atoms} (family names and their aliases, hash names
      with sizes, modes, paddings, curves, the SLH-DSA suffixes \code{s}/\code{f}, numbers and
      stop words such as \code{with}, \code{and}, \code{padding}, \code{encryption}). Atoms of at
      most two characters match only on word boundaries, so that no family is read out of an
      arbitrary word. Characters not covered by any atom form the \emph{residue}.
\item \textbf{Interpretation stage.} The primary family is chosen by a fixed priority
      (KDF $>$ DRBG $>$ MAC $>$ PQC $>$ classical public-key $>$ AEAD $>$ symmetric $>$ hash);
      numbers are attached as parameter sets according to the family's parameter-set kind;
      a hash atom becomes the hash slot; \code{chacha20}+\code{poly1305} becomes the AEAD family;
      two key-establishment families form a hybrid; pre-standard names (\code{kyber},
      \code{dilithium}, \code{falcon}, \code{sphincs+}) are mapped to the FIPS families with their
      parameter-set renaming (Dilithium 2/3/5 $\mapsto$ ML-DSA 44/65/87) and recorded in the notes.
\item \textbf{Hint stage.} Hints fill slots the name leaves at $\bot$; when a hint contradicts a
      name-derived slot the name wins and a \code{hint-conflict} note is recorded (finding C20).
\item \textbf{Convention stage.} Inadmissible slots are dropped; a padding on an AEAD or
      stream-like mode (gcm, ccm, ctr, ofb, cfb, xts, kw, \ldots) is dropped; an RSA signature
      scheme with a hash and no explicit padding is RSASSA-PKCS1-v1\_5
      (\code{rsa-pkcs1v15-}$h$); the primitive is derived from the family and the slots
      (a block cipher in an AEAD mode is \code{ae}; RSA with OAEP is \code{pke}, with PSS or a
      hash is \code{signature}; a hybrid is \code{combiner}).
\end{enumerate}
The confidence is $1$ when the residue is empty and no foreign family was ignored, $0.8$ when a
foreign family was ignored, $0.7$ when the residue is non-empty, and $0$ for $\bot$.
\end{definition}

\begin{theorem}[Fixed points of the grammar]\label{thm:fixpoint}
For every descriptor $\delta$ that respects $\mathrm{slots}(f)$ and whose family is in the
knowledge base, $\canon(\sem{\delta}, \emptyset) = (\delta, 1, \emptyset)$.
\end{theorem}
\begin{proof}[Proof sketch]
$\sem{\delta}$ contains only atoms of the lexicon separated by \code{-}; segmentation recovers the
atoms exactly (every slot value is itself an atom or a hash-atom plus number), the family is
unique, and the interpretation stage re-assigns each atom to the slot it came from because slot
kinds are disjoint (numbers to $\pi$, curve atoms to $\kappa$, etc.). The reference
implementation tests this property over its catalogue of names.
\end{proof}

\begin{remark}
Theorem~\ref{thm:fixpoint} is what makes identifiers \emph{stable}: re-canonicalising a
\twocbom{} document (for instance when it is itself fed back as an input) does not change any
identifier.
\end{remark}

\subsection{Cipher suites}
\begin{definition}[Cipher-suite decomposition]\label{def:suite}
An IANA cipher-suite name of the form \code{TLS_<kex>_<auth>_WITH_<cipher>_<hash>} decomposes into
constituents with roles: key exchange (\code{ECDHE}, \code{ECDH} $\mapsto$ \code{ecdh};
\code{DHE}, \code{DH} $\mapsto$ \code{dh}; \code{RSA} $\mapsto$ \code{rsa}; \code{PSK}),
authentication (\code{RSA}, \code{ECDSA}, \code{DSS} $\mapsto$ \code{dsa}; absent for
\code{anon}), the cipher $c$ (canonicalised by Definition~\ref{def:canon}, e.g.\
\code{3DES_EDE_CBC} $\mapsto$ \code{tdea-cbc}, \code{NULL} $\mapsto$ \code{null}) and the PRF
hash $h$ (\code{SHA} $\mapsto$ \code{sha1}, \code{SHA256} $\mapsto$ \code{sha2-256}, \ldots).
TLS~1.3 suites (\code{TLS_AES_256_GCM_SHA384}, no \code{WITH}) have no key-exchange
constituent; their \emph{kexKnown} flag is false and the key exchange \MUST{} be taken from the
protocol's references (\code{cryptoRefArray} groups). Hexadecimal identifiers are resolved
through the knowledge base (\code{0x1301} $\mapsto$ \code{TLS_AES_128_GCM_SHA256}).
\end{definition}

\subsection{Identity of the other asset kinds}
\begin{definition}[Identity functions and identity bases]\label{def:idfun}
\begin{itemize}[nosep]
\item \textbf{Certificate}, in order of preference ($H_{24}$ denotes the first 24 hex digits of
      SHA-256 and $x'$ the distinguished-name normalisation of $x$: attribute types upper-cased,
      spacing collapsed):
      \begin{itemize}[nosep]
      \item $\codem{2cbom:cert:fp:}\langle alg\rangle\codem{:}\langle hex\rangle$ when a fingerprint is available (basis \emph{fingerprint});
      \item $\codem{2cbom:cert:dn:}H_{24}(\JCS(\text{subject}', \text{issuer}', \text{serial}, \text{notBefore}, \text{notAfter}))$ (basis \emph{distinguished-name});
      \item $\codem{2cbom:cert:loc:}H_{24}(\text{location})$ (basis \emph{location}); otherwise \emph{unresolved}.
      \end{itemize}
\item \textbf{Key material} (with $\tau \in V_{\mathrm{mat}}$ the material type), in order of preference:
      \begin{itemize}[nosep]
      \item $\codem{2cbom:key:}\tau\codem{:fp:}\langle alg\rangle\codem{:}\langle hex\rangle$ by fingerprint;
      \item $\codem{2cbom:key:}\tau\codem{:id:}H_{24}(\text{identifier})$ by identifier;
      \item $\codem{2cbom:key:}\tau\codem{:loc:}H_{24}(\JCS(\text{algorithm}, \text{location}, \text{size}))$ by location;
      \item $\codem{2cbom:key:}\tau\codem{:src:}H_{24}(\text{source}, \text{ref})$ otherwise (basis \emph{source-local}).
      \end{itemize}
\item \textbf{Protocol.} $\id = \codem{2cbom:proto:}\tau\codem{-}v$ with $\tau \in V_{\mathrm{proto}}$ and $v$ the normalised version
      (\code{TLSv1.2}, \code{tls1.2}, \code{1.2} $\mapsto$ \code{1.2}; SSL versions are written
      \code{ssl3.0}).
\item \textbf{Component.} \code{2cbom:comp:} followed by the package URL without qualifiers,
      else the CPE, else \code{[group/]name[@version]} with normalised tokens.
\end{itemize}
Identity bases are totally ordered, fingerprint $>$ identifier $>$ distinguished-name $>$ location
$>$ source-local $>$ unresolved; the strongest available basis \MUST{} be used, and bases below
\emph{distinguished-name} raise finding C15.
\end{definition}

\subsection{Subsumption}
\begin{definition}[Specificity and subsumption]\label{def:subsume}
For descriptors of the same family, $\delta \sqsubseteq \delta'$ ($\delta'$ \emph{extends}
$\delta$) iff every slot defined in $\delta$ has the same value in $\delta'$ and $\delta'$ defines
at least one further slot. A raw detection $d$ with descriptor $\delta$ is \emph{subsumed} by a
raw detection $d'$ with descriptor $\delta' \sqsupset \delta$ iff $d$ and $d'$ carry evidence
with the same normalised location and the same line. A certificate detected by distinguished
name is subsumed by a certificate detected by fingerprint with the same name key. A subsumed
detection contributes its raw name (as alias), its evidence and its tool-reported attributes to
the subsuming asset but not its identity, and no asset is created for it. When several
detections subsume $d$, the most specific (most defined slots), then the lexicographically
smallest identifier, is chosen.
\end{definition}

% ==========================================================================================
\section{Quantum-security assessment model}\label{sec:assessment}
% ==========================================================================================
\begin{definition}[Assessment]\label{def:assessment}
An assessment is $\alpha = (q, \lambda_q, \lambda_c, \iota, \gamma, \varphi, \pi, \mathrm{rationale})$
with quantum class $q \in \Q$, NIST quantum security level $\lambda_q \in \{0,\dots,5\}\cup\{\bot\}$
(0: broken by Shor's algorithm; 1--5: the NIST PQC security categories, 1 $\equiv$ AES-128 key
search, 2 $\equiv$ SHA-256 collision, 3 $\equiv$ AES-192, 4 $\equiv$ SHA-384 collision,
5 $\equiv$ AES-256), classical security strength $\lambda_c \in \mathbb{N}\cup\{\bot\}$
(NIST~SP~800-57 bits), NIST~IR~8547 status $\iota \in V_{\mathrm{ir}}$, CNSA~2.0 membership
$\gamma \in \{0,1\}$, FIPS approval $\varphi \in \{0,1,\bot\}$, migration priority
$\pi \in V_{\mathrm{prio}}$ and a list of human-readable rationale strings, each produced by a
named rule.
\end{definition}

\begin{definition}[Order on quantum classes]
$\Q$ is totally ordered by \emph{badness}:
$$\text{classically-weak} \;<\; \text{quantum-vulnerable} \;<\; \text{unknown} \;<\; \text{quantum-safe},$$
and $\worst(X) = \min X$ for $X \subseteq \Q$, $\worst(\emptyset) = \text{unknown}$.
Placing \emph{unknown} below \emph{quantum-safe} is the central strictness choice of \twocbom:
an aggregate with an unassessed constituent is never reported safe.
\end{definition}

\begin{definition}[Level function]
The knowledge base defines $\lambda(f,\pi,\kappa,h) = (\lambda_c,\lambda_q)$ per family and
parameter set (SP~800-57 Table~2, FIPS~203/204/205, SP~800-208); for families whose strength is
inherited from a hash (HMAC, HKDF, PBKDF2, LMS/XMSS) $\lambda$ is that of the hash; for
elliptic-curve families it is that of the curve. The knowledge base is \emph{authoritative}:
tool-asserted levels that differ are kept as candidates and raise finding C13.
\end{definition}

\begin{definition}[Algorithm classification]\label{def:classify}
Let $P_{\mathrm{conf}} = \{\text{pke}, \text{kem}, \text{key-agree}\}$ and
$F_{\mathrm{conf}} = \{\text{encrypt}, \text{decrypt}, \text{encapsulate}, \text{decapsulate}, \text{keyderive}\}$
be the confidentiality primitives and functions, and $H_{\mathrm{weak}} = \{\text{sha1}, \text{md5},
\text{md4}, \text{md2}\}$. For an algorithm asset with descriptor $\delta$, crypto functions
\emph{fn} and family entry $F$ (whose class is quantum-safe for symmetric, hash and PQC
families, quantum-vulnerable for RSA, DSA, DH, ECDH, ECDSA, EdDSA, X25519/X448, SM2 and
ElGamal, and unknown for families absent from the knowledge base):
\begin{align*}
q &= \begin{cases}
  \text{classically-weak} & \text{if } F.\text{class} = \text{classically-weak} \text{ or } \lambda_c < 112 \text{ or } h \in H_{\mathrm{weak}}\\
  F.\text{class} & \text{otherwise}
\end{cases}\\
\iota &= \begin{cases}
  \text{acceptable} & q = \text{quantum-safe}\\
  \text{disallowed} & q = \text{classically-weak}, \text{ or } q = \text{quantum-vulnerable} \wedge \lambda_c < 112\\
  \text{deprecated-after-2030} & q = \text{quantum-vulnerable} \wedge \lambda_c = 112\\
  \text{disallowed-after-2035} & q = \text{quantum-vulnerable} \wedge (\lambda_c > 112 \vee \lambda_c = \bot)\\
  \text{unknown} & \text{otherwise}
\end{cases}\\
\pi &= \begin{cases}
  \text{immediate} & q = \text{classically-weak}\\
  \text{high} & q = \text{quantum-vulnerable} \wedge (\rho \in P_{\mathrm{conf}} \vee \text{fn} \cap F_{\mathrm{conf}} \ne \emptyset)\\
  \text{medium} & q = \text{quantum-vulnerable} \text{ otherwise}\\
  \text{none} & q = \text{quantum-safe} \wedge \gamma = 1\\
  \text{low} & q = \text{quantum-safe} \wedge \gamma = 0\\
  \text{unknown} & q = \text{unknown}.
\end{cases}
\end{align*}
$\gamma = 1$ iff $(f,\pi)$ is in the CNSA~2.0 set \{aes-256, sha2-384, sha2-512, ml-kem-1024,
ml-dsa-87, lms, hss, xmss, xmss-mt\}; $\varphi$ is the FIPS approval of the family restricted to
approved curves and non-broken hashes. The \emph{high} priority captures exposure to
harvest-now-decrypt-later attacks.
\end{definition}

\begin{definition}[Aggregation for composite assets]\label{def:aggregate}
Let $\mathrm{dep}(a)$ be the set of assets $a$ depends on for its security:
\begin{itemize}[nosep]
\item hybrid combiner: its constituents $\chi$; $q = $ classically-weak if any constituent is,
      quantum-safe if every PQ constituent is quantum-safe, unknown otherwise; $\lambda_q$ is the
      minimum over the PQ constituents, $\lambda_c$ the minimum over all;
\item cipher-suite-shaped algorithm, certificate (signature and subject-key algorithms), key
      material (the algorithms it instantiates), protocol (all constituents of all suites, all
      referenced assets including certificates): $q = \worst\{q_b : b \in \mathrm{dep}(a)\}$,
      $\lambda_q = \min$, $\lambda_c = \min$, $\gamma = \bigwedge$, $\pi = \worst$ of the
      constituent priorities (\emph{medium} for a quantum-vulnerable certificate, whose key
      authenticates rather than encrypts);
\item protocol version policy: SSL~2.0/3.0, TLS~1.0/1.1, SSH~1 and IKEv1 are classically-weak
      regardless of suites; a TLS~1.3 protocol whose suites have $\textit{kexKnown} = 0$ and
      that references no key-establishment asset is \emph{unknown};
\item symmetric key material without an algorithm: by size $n$, classically-weak if $n<112$,
      otherwise quantum-safe with $\lambda_q = 5,3,1$ for $n \ge 256, 192, 128$.
\end{itemize}
\end{definition}

\begin{proposition}[Monotonicity]
$\worst$ is monotone and idempotent; hence adding a constituent to a composite asset can only
lower its class, and the assessment of an asset is a function of the assessments of its
dependencies. Assessment therefore terminates on the dependency DAG (constituents are never
composites of their parents).
\end{proposition}

% ==========================================================================================
\section{Constraints and conformance profiles}\label{sec:constraints}
% ==========================================================================================
Let $D$ be a document, $\Assets_{\mathrm{alg}}$ the algorithm assets, etc., and let
$\mathrm{ev}(a) = E_a$, $\mathrm{rel}_\rho(a) = \{b : (\kappa_a,\kappa_b,\rho,\cdot,\cdot) \in \Rels\}$.

\setcounter{constraint}{0}
\begin{constraint}[Identity]\label{c:identity}
$\forall a,b \in \Assets \cup \Comps:\ \kappa_a = \kappa_b \Rightarrow a = b$, and in the strict
profile $\forall a \in \Assets_{\mathrm{alg}}: u_a = 0$. \hfill\emph{Finding C01}
\end{constraint}
\setcounter{constraint}{1}
\begin{constraint}[Evidence totality and location]
$\forall a \in \Assets: \mathrm{ev}(a) \neq \emptyset$ (C02), and $\exists e \in \mathrm{ev}(a)$
with $\ell_e \ne \bot$ (C02b, warning in strict, info in core).
\end{constraint}
\setcounter{constraint}{2}
\begin{constraint}[Assessment totality]
$\forall a \in \Assets: \alpha_a \ne \bot$, and in the strict profile $q_a \ne \text{unknown}$. \hfill\emph{C03}
\end{constraint}
\setcounter{constraint}{3}
\begin{constraint}[Referential integrity]
$\forall (\kappa_1,\kappa_2,\cdot,\cdot,\cdot) \in \Rels: \kappa_1,\kappa_2 \in \id(\Assets \cup \Comps)$,
and every \code{bom-ref} used by a source resolves within that source. \hfill\emph{C04}
\end{constraint}
\setcounter{constraint}{4}
\begin{constraint}[Certificate closure]
$\forall a \in \Assets_{\mathrm{cert}}: \mathrm{rel}_{\text{signed-by}}(a) \neq \emptyset$ (C05,
error in strict) and $\mathrm{rel}_{\text{subject-key-algorithm}}(a) \ne \emptyset$ (C06, warning).
\end{constraint}
\setcounter{constraint}{6}
\begin{constraint}[Protocol closure]
$\forall a \in \Assets_{\mathrm{proto}}$: $a$ has at least one cipher suite or referenced asset
(C07), every suite decomposes or references canonical algorithms (C08), and a TLS~1.3 protocol
records its key-exchange groups (C19).
\end{constraint}
\setcounter{constraint}{8}
\begin{constraint}[Key-material closure]
$\forall a \in \Assets_{\mathrm{mat}}$ with type $\in$ \{private-key, public-key\}:
$\mathrm{rel}_{\text{instantiates}}(a) \ne \emptyset$ (error in strict); for secret keys a
warning. \hfill\emph{C09}
\end{constraint}
\setcounter{constraint}{9}
\begin{constraint}[Parameterisation]\label{c:param}
$\forall a \in \Assets_{\mathrm{alg}}$ with family $f$: every slot in $\mathrm{requires}(f)$
is defined unless a slot in $\mathrm{requiresUnless}(f)$ is defined (an RSA signature scheme
identifier need not carry the modulus size). In the strict profile this is an error for
detections with located evidence and a warning for assets implied by cipher-suite names, whose
parameters are not part of the name. \hfill\emph{C10}
\end{constraint}
\setcounter{constraint}{10}
\begin{constraint}[Closed vocabularies]
Every enumerated attribute takes its value in the vocabulary of Definition~\ref{def:vocab}. \hfill\emph{C11}
\end{constraint}
\setcounter{constraint}{11}
\begin{constraint}[Temporal validity]
Certificate dates are ISO~8601 and $\text{notBefore} \le \text{notAfter}$. \hfill\emph{C12}
\end{constraint}
\setcounter{constraint}{12}
\begin{constraint}[Level consistency]
A tool-asserted $\lambda_c$ or $\lambda_q$ that differs from the knowledge base is reported
(warning) and overridden. \hfill\emph{C13}
\end{constraint}
\setcounter{constraint}{13}
\begin{constraint}[Lineage]
$\Srcs \ne \emptyset$; every source carries digest, format, spec version, tool and trust; every
attribute conflict is recorded (C14); pre-standard names are recorded (C16); non-standard spec
versions are recorded (C17); secret values are never carried (C18, Section~\ref{sec:security}).
\end{constraint}
\setcounter{constraint}{20}
\begin{constraint}[Determinism]
$D$ is serialised as in Section~\ref{sec:serialisation}; the digest verifies and the serial
number is derived from it. \hfill\emph{(validator)}
\end{constraint}

\begin{definition}[Profiles and conformance]
The \emph{core} profile demands C01 (first part), C02, C04, C11, C12, lineage and determinism as
errors and reports the remaining constraints as warnings or information. The \emph{strict}
profile, intended for PQC compliance, additionally raises C01 (second part), C03, C05, C09 (for
asymmetric keys) and C10 (for located detections) to errors. A document is \emph{conformant}
to a profile iff it contains no finding of severity \emph{error} under that profile;
$\Gamma.\text{conformant}$ records this.
\end{definition}

\begin{table}[ht]\centering\footnotesize
\begin{tabularx}{\textwidth}{lXll}
\toprule
Code & Constraint & strict & core\\
\midrule
C01-IDENTITY-UNRESOLVED & unresolved algorithm name & error & warning\\
C02-EVIDENCE-MISSING / -UNLOCATED & no evidence / no located evidence & error / warning & error / info\\
C03-ASSESSMENT-UNKNOWN & quantum class undetermined & error & warning\\
C04-DANGLING-RELATION & reference to unknown entity & warning & warning\\
C05-CERT-SIGNATURE-UNLINKED & certificate without signature algorithm & error & warning\\
C06-CERT-KEY-UNLINKED & certificate without subject-key algorithm & warning & warning\\
C07-PROTOCOL-SUITES-MISSING & protocol without suites or references & warning & warning\\
C08-SUITE-UNRESOLVED & cipher suite not decomposable & warning & warning\\
C09-KEY-ALGORITHM-UNLINKED & key material without algorithm & error$^{*}$ & warning\\
C10-PARAMETER-SET-MISSING & required slot undefined & error$^{\dagger}$ & warning\\
C11-VOCABULARY-VIOLATION & value outside closed vocabulary & error & error\\
C12-CERT-VALIDITY-INVALID & malformed or inverted validity & error & error\\
C13-LEVEL-INCONSISTENT & tool level contradicts catalogue & warning & warning\\
C14-SOURCE-CONFLICT & tools disagree, resolved by precedence & info & info\\
C15-WEAK-IDENTITY & identity by location / source-local & warning & warning\\
C16-PRESTANDARD-NAME & pre-standard name mapped & info & info\\
C17-SPEC-VERSION-NONSTANDARD & non-standard source spec version & info & info\\
C18-SECRET-VALUE-DROPPED & key material value discarded & warning & warning\\
C19-PROTOCOL-KEX-UNKNOWN & TLS~1.3 without recorded groups & warning & warning\\
C20-HINT-MISMATCH & structured field contradicts name & warning & warning\\
P01 / P02 & classically weak / quantum-vulnerable asset & warning / info & warning / info\\
I01 / I02 & subsumption / key-size refinement applied & info & info\\
\bottomrule
\end{tabularx}
\caption{Finding codes (prefix \texttt{2CBOM-}) and severities. $^{*}$asymmetric key types;
$^{\dagger}$detections with located evidence.}
\end{table}

% ==========================================================================================
\section{Canonical serialisation and integrity}\label{sec:serialisation}
% ==========================================================================================
\begin{definition}[Canonical JSON]
The JSON representation of $D$ uses the member names of Appendix~\ref{app:json}. Members whose
value is $\bot$ or an empty collection are omitted. All arrays are sorted: assets and
components by identifier, relations by $(\text{from}, \text{to}, \text{kind}, \text{via})$,
evidence by its key, findings by (severity rank, code, subject, message), aliases and sources
lexicographically. The canonical byte form is $\JCS(D)$, the RFC~8785 JSON Canonicalization
Scheme (sorted members, no insignificant whitespace, shortest number form, UTF-8); \twocbom{}
uses no non-integer numbers except evidence confidences, which \MUST{} be representable exactly
as decimal fractions with at most four digits.
\end{definition}

\begin{definition}[Digest scope, digest and serial number]
Let $D^{\circ}$ be $D$ without the members \code{integrity}, \code{serialNumber} and
\code{generated.timestamp}. Then
\begin{gather*}
\digest(D) = \codem{sha256:}\mathrm{hex}(\mathrm{SHA\text{-}256}(\JCS(D^{\circ}))),\\
\textit{sn}(D) = \codem{urn:uuid:}\mathrm{UUIDv5}(N_{\twocbom}, \digest(D)),
\end{gather*}
with the fixed namespace $N_{\twocbom} = \mathrm{UUIDv5}(\mathrm{NAMESPACE\_URL},
\codem{https://2cbom.org/ns/1.0})$. The integrity block records the algorithm
(\code{sha-256}), the canonicalisation (\code{RFC8785}), the digest and the scope.
\end{definition}

\begin{theorem}[Serialisation determinism]
If two documents have equal values of $D^{\circ}$ then they have equal digests and serial
numbers; conversely a validator recomputing $\digest(D)$ detects any modification of
$D^{\circ}$, including modifications of assessments. The timestamp is excluded so that
re-generation from the same inputs is byte-comparable.
\end{theorem}

% ==========================================================================================
\section{Merge semantics}\label{sec:merge}
% ==========================================================================================
\subsection{Assertion model}
\begin{definition}[Fact alphabet]
A \emph{fact} is an immutable, typed tuple. The fact types are: \textsc{Source}; the raw
facts \textsc{RawComponent}, \textsc{RawAsset}$(s, r, \tau, \nu, \text{fields})$,
\textsc{RawEvidence}, \textsc{RawRelation} obtained syntactically from a source; the derived
facts \textsc{Candidate}$(s, r, \kappa, \ldots)$ (a proposed identity), \textsc{Subsumed},
\textsc{Asset}$(\kappa,\tau)$, \textsc{Component}, \textsc{Member}$(s,r,\kappa)$ (raw-to-canonical
membership), \textsc{Alias}, \textsc{Attr}$(\kappa, n, v, o, \sigma, \tau_o, c)$ (an attribute
\emph{candidate} with value $v$, origin $o$, strength $\sigma$, trust $\tau_o$ and confidence
$c$), \textsc{Resolved}$(\kappa, n, v, \text{conflict})$, \textsc{Conflict}, \textsc{Evidence},
\textsc{Relation}, \textsc{Suite}, \textsc{Assessment} and \textsc{Finding}. A \emph{working
memory} $W$ is a finite set of facts.
\end{definition}

\begin{definition}[Lifting of sources]
For a source $s$, $\mathrm{facts}(s)$ is the set of \textsc{Source}, \textsc{Raw}$\ast$ facts
obtained by the format adapter of Section~\ref{sec:mapping}. The adapter is syntactic
(renaming, kebab-casing, flattening nested components, normalising locations); it takes no
semantic decision. A raw asset without any occurrence receives one \textsc{RawEvidence} with
method \emph{declared}.
\end{definition}

\subsection{Rules and the fixpoint}
\begin{definition}[Rule]
A rule is $R = (\text{name}, \text{salience} \in \mathbb{Z}, \text{when}, \text{then})$ where
$\text{when}(W)$ enumerates \emph{bindings} (tuples of facts and values for which $R$ applies)
and $\text{then}(W, \beta)$ returns a finite set of facts to assert. Rules are
\emph{monotone}: they only add facts. A rule \MAY{} test the absence of a fact only if every
rule that can produce such a fact has strictly higher salience (\emph{stratification}); the
single such rule in \twocbom{} is materialisation, which tests \textsc{Subsumed} facts produced
by the subsumption rules at salience 850.
\end{definition}

\begin{definition}[Recognise--act cycle]\label{def:cycle}
Given a rule base $\mathcal{R}$ ordered by $(-\text{salience}, \text{registration order})$ and a
set $\Phi$ of already fired (rule, binding) pairs, one cycle selects the first rule $R$ in that
order having a binding $\beta \in \text{when}(W)$ with $(R,\beta) \notin \Phi$, and fires
\emph{all} such bindings of $R$: $W \leftarrow W \cup \bigcup_\beta \text{then}(W,\beta)$,
$\Phi \leftarrow \Phi \cup \{(R,\beta)\}$ (\emph{refraction}). The run terminates when no such
rule exists. Every firing is appended to a trace $(\text{cycle}, R, \beta, \text{derived facts})$,
which is the explanation facility.
\end{definition}

\begin{lemma}[Termination and uniqueness]\label{lem:fixpoint}
Let $T_{\mathcal R}(W) = W \cup \bigcup_{R \in \mathcal R} \bigcup_{\beta \in \text{when}(W)} \text{then}(W,\beta)$.
If every \emph{then} is deterministic and depends on $W$ only through facts, if the number of
distinct facts derivable from a finite $W_0$ is finite, and if stratification holds, then the
recognise--act cycle started at $W_0$ terminates and its final state is $\lfp_{W_0}(T_{\mathcal R})$,
independently of the registration order of rules of equal salience.
\end{lemma}
\begin{proof}[Proof sketch]
Monotonicity and finiteness give termination (each cycle fires at least one new (rule, binding)
pair from a finite set). Any firing order reaches the same closure because \emph{then} is a
function of $W$ and bindings are never invalidated (no retraction); stratification ensures that
the single negation is evaluated on a complete stratum, which makes $T_{\mathcal R}$ monotone on
the fact lattice; the least fixpoint above $W_0$ is unique (Knaster--Tarski).
\end{proof}

\subsection{The \twocbom{} rule base}
The rule base is organised in salience bands; Table~\ref{tab:rules} lists it. Three rule groups
carry the semantic content of the merge:

\paragraph{Resolution of attributes (salience 500).}
\begin{definition}[Resolution order]\label{def:order}
For \textsc{Attr} candidates of one $(\kappa, n)$ define the total order
$$a \prec b \iff (\text{rank}(o_a), \sigma_a, \tau_a, c_a, v_a^{-1}, o_a^{-1}) >_{\mathrm{lex}} (\text{rank}(o_b), \sigma_b, \tau_b, c_b, v_b^{-1}, o_b^{-1})$$
with $\text{rank}(\text{kb}) = 3 > \text{rank}(\text{derived}) = 2 > \text{rank}(\text{tool}) = 1$
and $v^{-1}, o^{-1}$ denoting lexicographic comparison of the JCS encoding of the value and of
the origin identifier in \emph{ascending} order. \textsc{Resolved}$(\kappa,n,\cdot)$ takes the
$\prec$-minimum. Multi-valued attributes (crypto functions, certification levels, standards,
pre-standard names, cipher suites, references) are resolved by set union. A \textsc{Conflict} is
recorded when two tool origins assert different values, or when a knowledge-base or derived
value overrides a tool value.
\end{definition}
Strength $\sigma$ is the tool's declared specificity for the attribute (from the tool profile,
e.g.\ CBOMkit for mode, padding and parameter set; cbomkit-theia for certificate and key
attributes); trust $\tau$ is the tool profile's trust; the knowledge base is authoritative for
$\lambda_c, \lambda_q$, OID, standards, FIPS status and the family's quantum class; identity
slots are \emph{derived} from the canonical descriptor and therefore outrank tools.

\paragraph{Closure (salience 700).} Strings naming algorithms inside other assets
(\code{signatureAlgorithm}, \code{subjectPublicKeyAlgorithm}, \code{algorithm} of key material,
cipher-suite names, hybrid constituents) are canonicalised and materialised as assets with
\emph{derived} evidence that names the parent and the rule; references by \code{bom-ref} are
followed within the source. Key size refines a generic family (\code{rsa} with size 2048 also
instantiates \code{rsa-2048}, finding I02).

\paragraph{Subsumption (salience 850).} Definition~\ref{def:subsume}, applied before
materialisation (stratified).

\begin{table}[ht]\centering\footnotesize
\begin{tabularx}{\textwidth}{rXX}
\toprule
Salience & Rule & Effect\\
\midrule
900 & normalize.\{algorithm, certificate, material, protocol, component\} & \textsc{Candidate} identities (Section~\ref{sec:identity})\\
850 & resolve.subsume-by-location, resolve.subsume-certificate-by-dn & \textsc{Subsumed}\\
800 & resolve.materialize & \textsc{Asset}/\textsc{Component}, \textsc{Member}, \textsc{Alias}, identity \textsc{Attr}s\\
790 & resolve.attributes, resolve.component-attributes & provenance-tagged \textsc{Attr}s\\
785--775 & resolve.evidence, resolve.relations, resolve.dangling-relation & \textsc{Evidence}, typed \textsc{Relation}s, C04\\
700 & closure.\{certificate-signature, certificate-subject-key, material-algorithm, material-secured-by, protocol-suites, protocol-crypto-refs, constituents\} & referenced algorithms, \textsc{Suite}s, relations\\
690 & closure.refine-generic-by-material & sized instantiation\\
600, 590 & enrich.catalogue, enrich.hybrid-levels & knowledge-base \textsc{Attr}s\\
500 & settle.resolve-attribute & \textsc{Resolved}, \textsc{Conflict}\\
400--380 & assess.\{algorithm, certificate, material, protocol\} & \textsc{Assessment}\\
300--280 & check.C01 \ldots{} check.C20, check.P01-P02 & \textsc{Finding}\\
\bottomrule
\end{tabularx}
\caption{The \twocbom{} rule base (45 rules) of the reference implementation.}\label{tab:rules}
\end{table}

\subsection{The merge operator}
\begin{definition}[Merge]\label{def:merge}
For a finite set of sources $S = \{s_1,\dots,s_k\}$ and a profile $p$,
$$\merge_p(S) = \mathrm{assemble}\Big(\lfp_{W_0}(T_{\mathcal R})\Big),\qquad
  W_0 = \{\textsc{Config}(p)\} \cup \bigcup_{i} \mathrm{facts}(s_i),$$
where \emph{assemble} builds $D$ from the fixpoint (one asset per \textsc{Asset} fact with its
\textsc{Resolved} attributes, evidence, aliases, sources and assessment; relations grouped by
key; findings and conflicts collected) and seals it (Section~\ref{sec:serialisation}).
\end{definition}

\begin{theorem}[Order independence]\label{thm:merge}
$\merge_p$ is a function of the \emph{set} $S$: for any permutation of the sources and any
duplication of a source, $\merge_p$ yields the same $D^{\circ}$ and hence the same digest and
serial number. Equivalently, the binary operation $\merge_p(S_1 \cup S_2)$ is commutative,
associative and idempotent.
\end{theorem}
\begin{proof}
$W_0$ is a set union of per-source fact sets, each of which depends only on the source bytes
(source identifiers are content-derived, so a duplicated source contributes identical facts).
By Lemma~\ref{lem:fixpoint} the fixpoint is unique and order-free. The resolution order of
Definition~\ref{def:order} is total and refers only to fact contents (never to arrival order),
so every \textsc{Resolved} value is a function of its candidate set; set-union attributes are
order-free by construction. \emph{assemble} sorts every collection by content keys. Hence
$D^{\circ}$ is a function of $S$. Commutativity, associativity and idempotence follow from the
corresponding properties of set union.
\end{proof}

\begin{remark}[What is and is not decided by order]
The only place where ties are broken is Definition~\ref{def:order}; the tie-break is
lexicographic on value and origin identifier, both content-derived. Consequently two tools of
equal trust and specificity that disagree on a scalar attribute yield a deterministic choice
\emph{and} a conflict record, which is the auditable form of ``we could not tell''.
\end{remark}

\begin{figure}[ht]
\begin{lstlisting}
procedure MERGE(sources S, profile p):
    W  <- { Config(p) }
    for s in S:                               # order irrelevant (Theorem 3)
        W <- W U facts(s)                     # syntactic lifting by the format adapter
    Phi <- {}                                  # refraction set
    loop:
        pick the first rule R (by -salience, registration) with an unfired binding
        if none: break
        for every unfired binding b of R:     # fire all activations of R
            W <- W U then_R(W, b); Phi <- Phi U {(R, b)}; append (R, b, new facts) to TRACE
    D <- assemble(W);  D.integrity <- digest(scope(D));  D.serialNumber <- uuid5(N, D.integrity)
    return D, TRACE
\end{lstlisting}
\caption{The merge algorithm.}
\end{figure}

% ==========================================================================================
\section{Mappings from source formats}\label{sec:mapping}
% ==========================================================================================
\subsection{Format detection}
A JSON document is an IBM pre-standard CBOM iff \code{bomFormat} $=$ \code{"CBOM"} or
\code{specVersion} starts with \code{1.4-cbom}; a CycloneDX~1.4 document whose components
include type \code{crypto-asset} is treated the same way. A document with \code{bomFormat} $=$
\code{"CycloneDX"} and \code{specVersion} $\in \{1.5, 1.6, 1.7\}$ is read by the CycloneDX
adapter (1.5 and unknown versions with finding C17). Documents without markers are classified by
the presence of \code{cryptoProperties}.

\subsection{CycloneDX 1.6 and 1.7}
\begin{longtable}{p{5.6cm}p{5.6cm}p{3.8cm}}
\toprule
CycloneDX 1.6/1.7 & \twocbom{} & Note\\
\midrule
\endhead
\code{metadata.tools.components[0]} (or legacy \code{tools[]}) & source tool, version & trust from tool profile\\
\code{metadata.component} & root component & \\
\code{components[].type=cryptographic-asset} & asset & nested components flattened\\
\code{cryptoProperties.assetType} & $\tau_a$ & identity map\\
\code{algorithmProperties.{primitive, parameterSetIdentifier, curve, mode, padding}} & hints to $\canon$ & name wins on conflict (C20)\\
\code{algorithmProperties.{cryptoFunctions, executionEnvironment, implementationPlatform, certificationLevel, classicalSecurityLevel, nistQuantumSecurityLevel}} & tool attribute candidates & levels overridden by KB (C13)\\
\code{cryptoProperties.oid} & OID hint / candidate & \\
\code{certificateProperties.{subjectName, issuerName, notValidBefore, notValidAfter, certificateFormat, certificateExtension}} & certificate record & \\
\code{certificateProperties.signatureAlgorithmRef}, \code{subjectPublicKeyRef} & relations \code{signed-by}, \code{subject-key-algorithm} & followed within source\\
\code{relatedCryptoMaterialProperties.{type, id, state, algorithmRef, size, format, securedBy, dates}} & material record, \code{instantiates}, \code{secured-by} & \code{value} dropped (C18)\\
\code{protocolProperties.{type, version, cipherSuites[], cryptoRefArray}} & protocol record, \textsc{Suite}s, \code{negotiates}, \code{uses} & suites decomposed\\
\code{hashes[]} & fingerprint & identity basis\\
\code{evidence.occurrences[].{location, line, offset, symbol, additionalContext}} & evidence record & location normalised\\
\code{dependencies[].dependsOn} & \code{uses} (component$\to$asset), \code{depends-on} (component$\to$component) & \\
\code{dependencies[].provides} & \code{provides} & \\
\code{properties[]} & tool properties (kept verbatim) & \\
\bottomrule
\end{longtable}
CycloneDX~1.7 keeps the 1.6 cryptographic model; the adapter reads both with one code path and
passes unknown members through as tool properties.

\subsection{IBM pre-standard CBOM (\texttt{1.4-cbom-1.0})}
\begin{longtable}{p{5.6cm}p{5.6cm}p{3.8cm}}
\toprule
IBM 1.4-cbom-1.0 & \twocbom{} & Note\\
\midrule
\endhead
\code{bomFormat="CBOM"}, \code{specVersion="1.4-cbom-1.0"} & source format \code{ibm-cbom-json} & C17\\
\code{components[].type=crypto-asset} & asset & \\
\code{assetType=relatedCryptoMaterial} & \code{related-crypto-material} & \\
\code{algorithmProperties.variant} & name (when the component name is empty) & \\
\code{primitive} \code{blockcipher}, \code{streamcipher}, \code{keyAgree} & \code{block-cipher}, \code{stream-cipher}, \code{key-agree} & kebab-casing\\
\code{implementationLevel} \code{softwarePlainRam}, \ldots & \code{executionEnvironment} \code{software-plain-ram}, \ldots & \\
\code{relatedCryptoMaterialType} \code{privateKey}, \ldots & \code{type} \code{private-key}, \ldots & \\
\code{relatedCryptoMaterialProperties.algorithm} (string) + \code{size} & \code{instantiates} the canonicalised algorithm, size as parameter-set hint & I02\\
\code{certificateProperties.signatureAlgorithm}, \code{subjectPublicKeyAlgorithm} (strings) & canonicalised and linked (\code{signed-by}, \code{subject-key-algorithm}) & \\
\code{detectionContext[].{filePath, lineNumbers[], offsets[], keywords[], additionalContext}} & one evidence record per line & keyword as symbol\\
\code{confidenceLevel} & evidence confidence & \\
\code{dependencies[].dependencyType} \code{uses} / \code{implements} & relation kind & \\
\bottomrule
\end{longtable}

% ==========================================================================================
\section{Projection to CycloneDX 1.6}\label{sec:projection}
% ==========================================================================================
A \twocbom{} document projects onto a CycloneDX~1.6 BOM as follows: every component and asset
becomes a component with \code{bom-ref} equal to its canonical identifier; descriptor slots map to
\code{algorithmProperties} (modes and paddings outside the CycloneDX enumeration map to
\code{other} and are repeated under \code{properties} as \code{2cbom:algorithm:mode} /
\code{2cbom:algorithm:padding}); the hash slot, family, standards and constituents are carried as
\code{2cbom:algorithm:*} properties; assessments as \code{2cbom:assessment:*} properties;
relations \code{uses}, \code{depends-on} and \code{provides} map to \code{dependencies}, all other
kinds to \code{2cbom:relation:<kind>} properties of the metadata; evidence maps to
\code{evidence.occurrences} with tool and method prefixed to \code{additionalContext}; lineage,
profile, digest and conformance map to metadata properties. The projection is lossy with respect
to typed relations and provenance but \emph{lossless by annotation}: every \twocbom{} member is
recoverable from the properties. The projection's \code{serialNumber} is UUIDv5 of the
\twocbom{} digest.

% ==========================================================================================
\section{Security and privacy considerations}\label{sec:security}
% ==========================================================================================
\begin{itemize}[nosep]
\item A \twocbom{} document \MUSTNOT{} contain key material, passwords, tokens or other secret
      values; adapters discard \code{relatedCryptoMaterialProperties.value} and record C18.
      Fingerprints of public artefacts (certificates, public keys) are permitted.
\item The digest protects integrity but not authenticity; signatures over $\JCS(D^{\circ})$ are
      out of scope but straightforward (CycloneDX attestations or JWS over the canonical bytes).
\item Trust values order tool assertions; they are declared in the lineage and \MUST{} be reviewed
      when a tool is upgraded. The knowledge base version is recorded in the lineage for the same
      reason.
\item Evidence locations may reveal repository structure; documents intended for sharing \MAY{}
      strip prefixes (the implementation supports \code{--strip-prefix}) but \MUSTNOT{} remove
      evidence records, since C02 would then fail.
\item \emph{Unknown} assessments are reported as such and block strict conformance by design;
      consumers \MUSTNOT{} treat an unknown class as safe.
\end{itemize}

% ==========================================================================================
\appendix
\section{JSON representation}\label{app:json}
% ==========================================================================================
Member names (JSON Schema \code{2cbom-1.0.schema.json} of the reference implementation):
\begin{lstlisting}
{ "2cbomVersion": "1.0", "profile": "strict",
  "serialNumber": "urn:uuid:...", "generated": {"timestamp": "...", "generator": {...}},
  "integrity": {"algorithm": "sha-256", "canonicalization": "RFC8785", "digest": "sha256:...", "scope": "..."},
  "lineage": {"sources": [{"id": "src-36d615754603", "digest": "sha256:...", "format": "cyclonedx-json",
                           "specVersion": "1.6", "tool": {"name": "cbomkit", "version": "2.1.0"},
                           "trust": 80, "detectionMethod": "static-source", "assetCount": 6}],
              "ruleset": {"version": "1.0", "firedRules": {...}}, "knowledgeBase": {"version": "..."},
              "conflicts": [{"subject": "2cbom:alg:...", "attribute": "...", "chosen": ..., "candidates": [...], "rule": "..."}]},
  "components": [{"id": "2cbom:comp:pkg:maven/org.bouncycastle/bcprov-jdk18on@1.78", "type": "library", ...}],
  "assets": [{"id": "2cbom:alg:aes-128-gcm", "type": "algorithm", "name": "aes-128-gcm",
              "aliases": ["AES", "AES-128-GCM", "AES128-GCM"],
              "algorithm": {"family": "aes", "primitive": "ae", "parameterSet": "128", "mode": "gcm",
                            "cryptoFunctions": ["decrypt", "encrypt"], "classicalSecurityLevel": 128,
                            "nistQuantumSecurityLevel": 1, "oid": "2.16.840.1.101.3.4.1.6",
                            "standards": ["FIPS 197", "NIST SP 800-38D"], "executionEnvironment": "software-plain-ram"},
              "assessment": {"quantumClass": "quantum-safe", "nistQuantumSecurityLevel": 1,
                             "classicalSecurityLevel": 128, "nistIr8547": "acceptable", "cnsa2": false,
                             "fipsApproved": true, "migrationPriority": "low", "rationale": [...]},
              "evidence": [{"source": "src-36d615754603", "tool": "cbomkit", "method": "static-source",
                            "location": "src/main/java/com/example/payments/CryptoUtil.java", "line": 42,
                            "offset": 23, "context": "Cipher.getInstance(\"AES/GCM/NoPadding\")", "confidence": 1.0}],
              "sources": ["src-36d615754603", "src-6b22f6d672d9", "src-e5fc2a359bba"], "unresolved": false}],
  "relations": [{"from": "2cbom:proto:tls-1.2", "to": "2cbom:alg:tdea-cbc", "kind": "negotiates",
                 "via": "TLS_RSA_WITH_3DES_EDE_CBC_SHA", "sources": ["src-6b22f6d672d9"]}],
  "findings": [{"code": "2CBOM-P01-CLASSICALLY-WEAK", "severity": "warning", "subject": "2cbom:alg:sha1",
                "message": "...", "rule": "check.P01-P02-readiness"}],
  "conformance": {"profile": "strict", "conformant": true, "errors": 0, "warnings": 16, "infos": 21},
  "summary": {"assetCount": 27, "componentCount": 7, "relationCount": 59,
              "assetsByQuantumClass": {...}, "assetsByMigrationPriority": {...}} }
\end{lstlisting}

\section{Knowledge-base excerpt}
The algorithm catalogue records, per family: primitive, admissible and required slots, parameter
sets with $(\lambda_c,\lambda_q)$, quantum class, governing standards, FIPS approval, CNSA~2.0
parameter sets, OIDs and pre-standard names. Examples:
\begin{center}\footnotesize
\begin{tabularx}{\textwidth}{lXlllX}
\toprule
Family & Parameter set & $\lambda_c$ & $\lambda_q$ & Class & Standard\\
\midrule
aes & 128 / 192 / 256 & 128 / 192 / 256 & 1 / 3 / 5 & quantum-safe & FIPS 197\\
sha2 & 256 / 384 / 512 & 128 / 192 / 256 & 2 / 4 / 5 & quantum-safe & FIPS 180-4\\
rsa & 2048 / 3072 / 4096 & 112 / 128 / 128$^{\ddagger}$ & 0 & quantum-vulnerable & FIPS 186-5\\
ecdsa, ecdh & p256 / p384 / p521 & 128 / 192 / 256 & 0 & quantum-vulnerable & FIPS 186-5, SP 800-186\\
ml-kem & 512 / 768 / 1024 & 128 / 192 / 256 & 1 / 3 / 5 & quantum-safe & FIPS 203\\
ml-dsa & 44 / 65 / 87 & 128 / 192 / 256 & 2 / 3 / 5 & quantum-safe & FIPS 204\\
slh-dsa & sha2-128s \ldots{} shake-256f & 128 / 192 / 256 & 1 / 3 / 5 & quantum-safe & FIPS 205\\
lms, hss, xmss & (hash) & 128 & 2 & quantum-safe & SP 800-208\\
tdea & 168 & 112 & 0 & classically-weak & SP 800-67 (withdrawn)\\
sha1 & -- & 80 & 0 & classically-weak & SP 800-131A\\
\bottomrule
\end{tabularx}
\end{center}
$^{\ddagger}$conservative (SP~800-57 lists 3072 $\to$ 128 and 7680 $\to$ 192). Catalogue entries are
data, versioned in the lineage; a compliance programme \SHOULD{} review them against the
standards in force.

\section{Reference implementation}
\code{twocbom} (Python~3.10+, no runtime dependencies) implements this specification:
\code{twocbom.canonical} (Sections~\ref{sec:identity} and~\ref{sec:serialisation}),
\code{twocbom.ingest} (Section~\ref{sec:mapping}), \code{twocbom.engine} and \code{twocbom.rules}
(Section~\ref{sec:merge}, Table~\ref{tab:rules}), \code{twocbom.validate} (Section~\ref{sec:constraints})
and \code{twocbom.export} (Section~\ref{sec:projection}). Its test suite checks
Theorem~\ref{thm:fixpoint} over the catalogue of names and Theorem~\ref{thm:merge} over
permutations and duplications of five source CBOMs (CBOMkit, cdxgen, cryptobom-forge, an IBM
pre-standard scanner and a CycloneDX~1.7 scanner), and the command \code{twocbom explain} replays
the trace of Definition~\ref{def:cycle} for any identifier.

\end{document}
